\documentclass[conference]{IEEEtran}
\IEEEoverridecommandlockouts

\usepackage{cite}
\usepackage{amsmath,amssymb,amsfonts}
\usepackage{algorithmic}
\usepackage{graphicx}
\usepackage{textcomp}
\usepackage{xcolor}
\usepackage{booktabs}
\usepackage{hyperref}
\usepackage{multirow}
\usepackage{array}

\def\BibTeX{{\rm B\kern-.05em{\sc i\kern-.025em b}\kern-.08em
    T\kern-.1667em\lower.7ex\hbox{E}\kern-.125emX}}

\begin{document}

\title{Smartphone Camera-Based Digital Biomarker Collection for AI-Driven Disease Prediction: A Systematic Review and Tiered Evaluation Framework}

\author{\IEEEauthorblockN{Anonymous Authors}
\IEEEauthorblockA{(Submitted for peer review)}}

\maketitle

\begin{abstract}
The proliferation of smartphones equipped with high-resolution cameras has created an unprecedented opportunity to collect digital biomarkers for non-invasive health monitoring and disease prediction. This paper presents the first systematic review and tiered evaluation framework for smartphone camera-based digital biomarkers across six major technology domains: remote photoplethysmography (rPPG), contact-based smartphone PPG, facial and skin image analysis, eye and gaze tracking, motion and gait analysis, and mental health assessment through facial expression analysis. We conducted a comprehensive literature search across PubMed, IEEE Xplore, Google Scholar, and other databases, identifying 42 studies published between 2019 and 2026, of which 30 were fully verified through a reference hallucination guard protocol. We propose a five-dimensional evaluation matrix assessing technical maturity, clinical validity, practicality, data availability, and regulatory readiness, yielding a three-tier prioritization framework. Tier~1 (immediately deployable) biomarkers include rPPG-based heart rate estimation (MAE 0.5--3 bpm), smartphone PPG atrial fibrillation detection (accuracy 98.5\%, FDA-cleared), diabetic retinopathy fundus screening (sensitivity 0.93, meta-analysis of 82 studies), neonatal jaundice detection (sensitivity 100\%, JAMA-validated), and rPPG-based SpO$_2$ estimation (MAE 1.274\%). Tier~2 encompasses 13 promising biomarkers requiring further validation, while Tier~3 includes 11 exploratory markers. We identify a critical research gap in women's health, specifically endometriosis and polycystic ovary syndrome (PCOS), where no camera-based biomarker studies exist despite strong physiological rationale linking rPPG heart rate variability to autonomic nervous system dysregulation in these conditions. We conclude with a translational roadmap integrating technical maturity, clinical validation needs, and practical deployment considerations.
\end{abstract}

\begin{IEEEkeywords}
digital biomarker, smartphone camera, remote photoplethysmography, rPPG, disease prediction, deep learning, mHealth, computer vision
\end{IEEEkeywords}

%==============================================================================
\section{Introduction}
%==============================================================================

The global smartphone penetration, exceeding 6.9 billion users in 2025, has positioned these ubiquitous devices as potentially transformative tools for decentralized health monitoring. Modern smartphones integrate high-resolution cameras (12--200 megapixels), powerful neural processing units (NPUs), and advanced computational photography pipelines, making them capable sensor platforms for extracting physiological signals without dedicated medical equipment.

Digital biomarkers---defined as objective, quantifiable physiological and behavioral measures collected by digital devices---have emerged as a bridge between consumer electronics and clinical medicine~\cite{dawadi2025scoping}. Among the various sensor modalities available on smartphones, the camera stands out for its ability to capture rich visual and photoplethysmographic data non-invasively.

Remote photoplethysmography (rPPG) enables contactless measurement of heart rate, blood pressure, and blood oxygen saturation by detecting subtle color changes on facial skin~\cite{debnath2025comprehensive}. Simultaneously, advances in computer vision and deep learning have enabled the detection of conditions ranging from skin cancer~\cite{naqvi2023skincancer} to Parkinson's disease~\cite{multimodalpd2025} from smartphone-captured images and videos.

Despite rapid progress, the field lacks a unified framework for evaluating and prioritizing these diverse biomarkers. Individual studies report impressive accuracy metrics, yet the translation from laboratory demonstrations to clinically validated, regulatory-approved tools remains challenging. Key barriers include variability in illumination conditions~\cite{acharya2025reliability}, skin tone bias~\cite{fibricheck2025}, motion artifacts, and the scarcity of large-scale, multi-ethnic validation studies.

This paper addresses three critical gaps in the literature:

\begin{enumerate}
\item \textbf{Fragmented knowledge}: No prior review has systematically cataloged smartphone camera-based biomarkers across all technology domains with a uniform evaluation framework.
\item \textbf{Lack of prioritization}: Existing reviews focus on single domains (e.g., rPPG or dermatology) without cross-domain comparison or deployment readiness assessment.
\item \textbf{Unexplored clinical applications}: We identify women's health---specifically endometriosis and polycystic ovary syndrome (PCOS)---as a significant research void where camera-based biomarkers could have substantial clinical impact.
\end{enumerate}

Our contributions are: (1) the first comprehensive taxonomy of smartphone camera-based digital biomarkers spanning six technology domains and seven disease categories; (2) a five-dimensional tiered evaluation framework integrating technical maturity, clinical validity, practicality, data availability, and regulatory readiness; (3) a translational roadmap from research prototype to clinical deployment; and (4) identification and preliminary analysis of women's health as a critical research frontier.

The remainder of this paper is organized as follows. Section~II reviews related work across rPPG, facial analysis, eye tracking, and motion analysis. Section~III describes our systematic review methodology including the reference hallucination guard protocol. Section~IV presents the biomarker taxonomy and tiered evaluation results. Section~V discusses technical challenges, clinical translation gaps, and women's health opportunities. Section~VI concludes with future directions.

%==============================================================================
\section{Related Work}
%==============================================================================

\subsection{Remote Photoplethysmography (rPPG)}

rPPG extracts cardiovascular signals from facial video by detecting micro-variations in skin color caused by blood volume pulse. Early approaches relied on signal processing techniques such as Independent Component Analysis (ICA), Plane-Orthogonal-to-Skin (POS), and Chrominance-based (CHROM) methods. Luo et al.~\cite{luo2019smartphone} demonstrated Transdermal Optical Imaging in a landmark study with 1,328 participants, achieving blood pressure estimation accuracy of 95.3\%/96.4\% for SBP/DBP within $\pm$5~mmHg.

The deep learning era has fundamentally transformed rPPG. Debnath and Kim~\cite{debnath2025comprehensive} provide a comprehensive review documenting the evolution from CNN-based architectures to Transformer models, with heart rate estimation achieving MAE of 0.5--3~bpm across benchmark datasets. Cheng et al.~\cite{cheng2024contactless} extended rPPG to contactless SpO$_2$ estimation using Spatial-Temporal Maps (STMap) with CNN, achieving MAE of 1.274\% and RMSE of 1.710\%, surpassing the international 4\% error threshold. A dedicated review of deep learning methods for rPPG~\cite{rppgreview2024frontiers} further cataloged advances in heart rate variability (HRV) extraction, reporting correlations of $r=0.85$--$0.95$ against ECG reference.

For clinical deployment, ReViSe~\cite{revise2022} demonstrated simultaneous real-time estimation of heart rate, respiratory rate, and SpO$_2$ within a smartphone application. More recently, Acharya et al.~\cite{acharya2025reliability} systematically evaluated rPPG reliability under low illumination and elevated heart rate, identifying conditions where current algorithms fail---a critical consideration for real-world deployment.

\subsection{Smartphone PPG for Cardiac Arrhythmia}

Contact-based smartphone PPG, where users place a finger over the rear camera, has achieved regulatory-grade performance for atrial fibrillation (AF) detection. FibriCheck~\cite{fibricheck2025} received FDA clearance following a multi-center validation across 236 participants using 10 different smartphone models, achieving accuracy of 98.5\%, sensitivity of 96.3\%, and specificity of 99.3\%. Gruwez et al.~\cite{gruwez2024realworld} independently validated smartphone PPG for AF in a real-world setting with 3,407 measurements, reporting sensitivity of 98.3\% and specificity of 99.9\%.

\subsection{Facial and Skin Image Analysis}

Computer vision applied to facial and skin images has enabled detection of diverse conditions. For anemia, Zhao et al.~\cite{zhao2024emoglobin} developed the eMoglobin application that analyzes conjunctival pallor from RAW smartphone images, achieving AUC of 0.92 for severe anemia (Hb$<$7~g/dL) in 426 emergency department patients. An exploratory study~\cite{bulbarconjunctiva2026} extended this to continuous blood count estimation from bulbar conjunctiva video.

In neonatal care, the BiliSG application~\cite{ngeow2024bilisg} detects jaundice using the Kramer principle applied to smartphone photos, achieving 100\% sensitivity and 70\% specificity (AUC 0.89) in a multi-ethnic cohort of 546 neonates published in JAMA Network Open.

For dermatology, comprehensive reviews~\cite{naqvi2023skincancer} document DenseNet169 achieving 92.25\% accuracy (F1=0.932) for skin cancer classification. Soenksen et al.~\cite{soenksen2021melanoma} at MIT demonstrated wide-field deep CNN for melanoma detection from whole-body photographs, enabling early screening outside clinical settings.

A meta-analysis of 82 studies encompassing 887,244 fundus images~\cite{drreview2025} established deep learning-based diabetic retinopathy screening at sensitivity 0.93 and specificity 0.90, with smartphone fundus photography enabling deployment in resource-limited settings.

For obstructive sleep apnea (OSA), craniofacial photography combined with questionnaires and CNN achieved sensitivity of 91.1\% and specificity of 79.2\% in 748 participants validated against polysomnography~\cite{osa2025jcsm}.

\subsection{Eye and Gaze Tracking}

Smartphone and tablet-based eye tracking has emerged as a biomarker for neurocognitive conditions. The m-ETA system~\cite{meta2025alzheimer} demonstrated community-based mild cognitive impairment (MCI) screening using tablet eye tracking. AI-driven eye tracking achieved AUC of 0.85 for distinguishing Alzheimer's disease (166 AD vs.\ 107 normal controls)~\cite{eyetrackingad2024}, while the VECA system~\cite{veca2024} combined VR-based eye tracking with machine learning for dementia screening.

For ADHD, tablet-based eye tracking integrated with continuous performance tasks achieved AUC of 0.965 with accuracy of 0.908 in 437 children~\cite{adhdeyetracking2024,adhdbiomarker2024}. Smartphone-based pupillometry using near-infrared cameras has also been explored as a pilot for neurological screening~\cite{pupilapp2022chi}.

Smartphone fundus photography has been applied to ophthalmological screening, including cataract detection using low-cost devices~\cite{cataract2024}.

\subsection{Motion and Gait Analysis}

Video-based gait and movement analysis enables detection of movement disorders. Sathya Bama and Bevish Jinila~\cite{sathyabama2022gait} applied Kernel PCA for vision-based Parkinson's disease (PD) identification from gait video, achieving 93.6\% accuracy. A multimodal smartphone assessment combining speech, hand movement, and gait data from 496 participants achieved integrated AUC of 0.86 for PD detection~\cite{multimodalpd2025}. CMSA-Net~\cite{cmsanet2025} proposed bilateral gait camera sensor fusion for portable PD screening.

For tremor assessment, computer vision using Mediapipe landmarks demonstrated significant correlation with clinical TETRAS scores in 66 essential tremor patients~\cite{tremorcv2024}. Privacy-preserving AI classified seven gait types from 743 mobile videos~\cite{gaitclassification2025}, demonstrating that skeleton-based features can replace raw video to protect patient privacy.

\subsection{Mental Health Assessment}

Camera-based mental health biomarkers span from controlled laboratory settings to in-the-wild deployment. Nepal et al.~\cite{nepal2024moodcapture} developed MoodCapture, a pioneering system that passively captures front-camera images during daily smartphone use, achieving balanced accuracy of 0.60--0.61 for major depressive disorder detection across 177 participants over 90 days. While modest in accuracy, MoodCapture represents the first in-the-wild depression detection system.

In controlled settings, Emoface~\cite{emoface2025} differentiated major depressive disorder from bipolar disorder using facial expression analysis of emotional stimulus videos, achieving 95.29\% accuracy in 700 participants. Fontes et al.~\cite{fontes2024stress} applied 1D-CNN to rPPG signals for stress detection, achieving 95.83\% on the UBFC-Phys benchmark. Additional work includes 16-class facial emotion recognition from smartphone videos~\cite{facialemotion2025} and audio-video fusion for depression severity classification~\cite{multimodaldepression2025}.

%==============================================================================
\section{Methodology: Systematic Review Protocol}
%==============================================================================

\subsection{Search Strategy}

We conducted a systematic literature search across PubMed, PubMed Central (PMC), IEEE Xplore, Google Scholar, arXiv, Semantic Scholar, Nature journals, JMIR, Springer, MDPI, Wiley, and JAMA Network from January 2015 through April 2026. Nineteen search queries were constructed combining technology terms (``remote photoplethysmography,'' ``smartphone camera,'' ``deep learning'') with clinical targets (``heart rate,'' ``blood pressure,'' ``anemia,'' ``skin cancer,'' ``Parkinson,'' ``depression'') and modality descriptors (``facial video,'' ``eye tracking,'' ``gait analysis'').

\subsection{Inclusion and Exclusion Criteria}

Studies were included if they: (1) utilized a smartphone camera or standard RGB camera as the primary sensor; (2) targeted disease prediction, detection, or physiological signal estimation; and (3) were published after 2015. We excluded studies using dedicated medical imaging equipment (CT, MRI, ultrasound) as the sole modality, wearable-only PPG devices (patches, rings), and purely clinical imaging systems without consumer-device applicability.

\subsection{Reference Hallucination Guard Protocol}

Given the known tendency of AI-assisted literature synthesis to generate plausible but non-existent citations, we implemented a rigorous reference hallucination guard protocol. Every cited paper was independently verified through:

\begin{enumerate}
\item \textbf{Direct DOI/URL access}: Papers with DOI or URL were fetched directly via PubMed Central, PLOS, JAMA Network Open, and publisher websites (10 papers verified via direct access).
\item \textbf{Cross-reference search}: Title, author, journal, and year combinations were searched on Google Scholar and PubMed to confirm existence and metadata consistency (32 papers verified).
\item \textbf{Content verification}: Where full text was accessible, key claims (sample size, model architecture, performance metrics) were compared against the source document.
\end{enumerate}

Each reference was classified as: \textbf{Verified} ($\checkmark$, 30 papers, 71.4\%), \textbf{Partially verified} (8 papers, 19.0\%, paper exists but with minor metric discrepancies), \textbf{Unconfirmed} (2 papers, 4.8\%), or \textbf{Suspected hallucination} ($\times$, 2 papers, 4.8\%). The two hallucinated references---a purported Vision Transformer-based anemia detection study (actual paper used VGG16+ResNet-50+InceptionV3 ensemble) and a purported craniofacial OSA meta-analysis (actual paper addressed wearable AI for sleep apnea)---were excluded from all analyses.

\subsection{Five-Dimensional Evaluation Framework}

Each identified biomarker was scored on five dimensions using a 1--5 Likert scale:

\begin{itemize}
\item \textbf{Technical Maturity} (TM): From proof-of-concept (1) to commercially deployed product (5).
\item \textbf{Clinical Validity} (CV): From case reports (1) to multi-center validation or meta-analysis with AUC$>$0.90 (5).
\item \textbf{Practicality} (PR): From requiring specialized equipment (1) to immediate collection on any smartphone (5).
\item \textbf{Data Availability} (DA): From proprietary/difficult to collect (1) to multiple open datasets available (5).
\item \textbf{Regulatory Readiness} (RR): From unclear regulatory pathway (1) to FDA/CE marking obtained (5).
\end{itemize}

The sum score (range 5--25) determined tier assignment: Tier~1 (20--25), Tier~2 (13--19), and Tier~3 (7--12).

%==============================================================================
\section{Results: Biomarker Taxonomy and Tiered Evaluation}
%==============================================================================

\subsection{Biomarker Taxonomy}

Our systematic review identified 33 distinct biomarkers from 40 verified studies (30 fully verified + 8 partially verified, excluding 2 hallucinated references), spanning six technology domains and seven disease categories. Table~\ref{tab:taxonomy} presents the taxonomy organized by technology type.

\begin{table}[t]
\caption{Biomarker Taxonomy by Technology Domain}
\label{tab:taxonomy}
\centering
\small
\begin{tabular}{p{2.2cm}p{2.8cm}cc}
\toprule
\textbf{Technology} & \textbf{Biomarkers} & \textbf{Studies} & \textbf{Algorithm} \\
\midrule
rPPG (contactless) & HR, HRV, BP, SpO$_2$, stress & 12 & CNN, Transformer \\
Smartphone PPG (contact) & AF detection & 2 & FibriCheck AI \\
Facial/skin imaging & Anemia, jaundice, skin cancer, DR, OSA & 9 & DenseNet, ViT, DCNN \\
Eye/gaze tracking & MCI/AD, ADHD, glaucoma, cataract & 8 & AI eye-tracking, DL \\
Motion/gait & PD, tremor, gait disorders & 6 & KPCA, Mediapipe, ML \\
Mental health & Depression, bipolar, stress & 6 & EfficientNet, 1D-CNN \\
\bottomrule
\end{tabular}
\end{table}

The disease category distribution reveals cardiovascular applications as the most mature (8 biomarkers; Tier~1: 3), followed by hematological/metabolic conditions (6 biomarkers; Tier~1: 2) and ophthalmology (3 biomarkers; Tier~1: 1). Neurological/movement (5 biomarkers) and mental health (6 biomarkers) domains remain predominantly in Tier~2 and Tier~3, reflecting their earlier developmental stage.

\subsection{Tier 1: Validated Biomarkers (Score 20--25)}

Five biomarkers met the threshold for immediate deployability, as detailed in Table~\ref{tab:tier1}. These share common characteristics: multi-center or meta-analytic validation, sample sizes exceeding 200, and either regulatory clearance or performance meeting established clinical standards.

\begin{table*}[t]
\caption{Tier 1: Validated Biomarkers Ready for Deployment}
\label{tab:tier1}
\centering
\small
\begin{tabular}{p{2.8cm}p{2.2cm}p{3.5cm}ccccccc}
\toprule
\textbf{Biomarker} & \textbf{Condition} & \textbf{Best Performance} & \textbf{TM} & \textbf{CV} & \textbf{PR} & \textbf{DA} & \textbf{RR} & \textbf{Total} \\
\midrule
rPPG Heart Rate & Cardiovascular & MAE 0.5--3 bpm & 5 & 5 & 5 & 5 & 3 & \textbf{23} \\
Smartphone PPG AF & Atrial fibrillation & Acc 98.5\%, Sen 96.3\% & 5 & 5 & 4 & 3 & 5 & \textbf{22} \\
DR Fundus Screening & Diabetic retinopathy & Sen 0.93, Spe 0.90 & 5 & 5 & 3 & 4 & 4 & \textbf{21} \\
Neonatal Jaundice (BiliSG) & Neonatal jaundice & Sen 100\%, Spe 70\% & 4 & 5 & 5 & 3 & 4 & \textbf{21} \\
rPPG SpO$_2$ & Hypoxemia & MAE 1.274\%, RMSE 1.710\% & 4 & 4 & 5 & 4 & 3 & \textbf{20} \\
\bottomrule
\end{tabular}
\end{table*}

\textbf{rPPG Heart Rate} (score 23) achieved the highest rating due to extensive validation across multiple datasets (UBFC-rPPG, PURE, VIPL-HR), comprehensive deep learning reviews~\cite{debnath2025comprehensive}, and demonstrated feasibility in both laboratory and in-the-wild conditions. The technology has matured from signal-processing methods to Transformer architectures that handle diverse lighting and motion conditions.

\textbf{Smartphone PPG AF Detection} (score 22) uniquely benefits from FDA clearance (FibriCheck~\cite{fibricheck2025}), multi-center validation across 10 smartphone models, and independent real-world confirmation~\cite{gruwez2024realworld}. The regulatory readiness score of 5 reflects the established SaMD (Software as a Medical Device) pathway.

\textbf{DR Fundus Screening} (score 21) leverages a meta-analysis of 82 studies encompassing 887,244 fundus images~\cite{drreview2025}, providing the highest evidence level in our review. The practicality score of 3 reflects the need for smartphone fundus adapters, a limitation compared to contactless approaches.

\textbf{BiliSG Neonatal Jaundice} (score 21) demonstrated exceptional clinical validity in a JAMA-published multi-ethnic study~\cite{ngeow2024bilisg}, with perfect sensitivity ensuring no missed cases---critical for a condition where delayed detection can cause kernicterus.

\textbf{rPPG SpO$_2$} (score 20) reached international accuracy thresholds ($<$4\% MAE)~\cite{cheng2024contactless}, with the highest practicality score (5) reflecting fully contactless measurement requiring only a smartphone front camera.

\subsection{Tier 2: Promising Biomarkers (Score 13--19)}

Thirteen biomarkers qualified as promising but requiring additional validation. Table~\ref{tab:tier2} presents selected results.

\begin{table}[t]
\caption{Tier 2: Promising Biomarkers (Selected)}
\label{tab:tier2}
\centering
\small
\begin{tabular}{p{2.5cm}p{2.5cm}c}
\toprule
\textbf{Biomarker} & \textbf{Performance} & \textbf{Score} \\
\midrule
rPPG Blood Pressure & SBP/DBP 95.3\%/96.4\% & 19 \\
rPPG HRV & ECG corr.\ $r$=0.85--0.95 & 18 \\
Anemia (eMoglobin) & AUC 0.92, $n$=426 & 18 \\
Skin cancer & DenseNet 92.25\% & 17 \\
PD gait multimodal & AUC 0.86, $n$=496 & 17 \\
MCI/AD eye tracking & AUC 0.85 & 16 \\
ADHD eye tracking & AUC 0.965, $n$=437 & 16 \\
Depression (Emoface) & Acc 95.29\%, $n$=700 & 15 \\
Tremor video analysis & TETRAS corr., $n$=66 & 15 \\
OSA craniofacial & Sen 91.1\%, $n$=748 & 14 \\
\bottomrule
\end{tabular}
\end{table}

Notable patterns emerge within Tier~2. \textbf{rPPG Blood Pressure} (score 19) represents the most clinically impactful biomarker pending multi-center validation; Luo et al.~\cite{luo2019smartphone} demonstrated feasibility in 1,328 participants, while video-based approaches~\cite{videobp2024} show promise for beat-by-beat monitoring. \textbf{ADHD eye tracking} achieves an unusually high AUC of 0.965~\cite{adhdbiomarker2024} but scores lower due to limited regulatory precedent and tablet-dependent hardware requirements. \textbf{Emoface}~\cite{emoface2025}, while achieving 95.29\% accuracy in differentiating depression from bipolar disorder, was evaluated in a controlled laboratory setting with curated emotional stimuli, limiting generalizability.

\subsection{Tier 3: Exploratory Biomarkers (Score 7--12)}

Eleven biomarkers remain at exploratory stages. MoodCapture~\cite{nepal2024moodcapture} (score 12) is notable as the only in-the-wild depression detection system, though its balanced accuracy of 0.60--0.61 is below clinical utility thresholds. rPPG-based stress detection~\cite{fontes2024stress} achieves 95.83\% on the UBFC-Phys benchmark but has not been validated in naturalistic settings. Privacy-preserving gait classification~\cite{gaitclassification2025} demonstrates an important paradigm for handling sensitive video data but requires further clinical correlation. Emerging work on noninvasive blood counts from conjunctiva video~\cite{bulbarconjunctiva2026} represents a frontier with transformative potential but remains at proof-of-concept stage.

%==============================================================================
\section{Discussion}
%==============================================================================

\subsection{Technical Challenges and the Architecture Transition}

Our review reveals a clear architectural transition across the 2019--2026 timeline. Early systems (2019--2022) relied predominantly on CNN architectures (ResNet, DenseNet, MobileNet), exemplified by DenseNet169 achieving 92.25\% for skin cancer~\cite{naqvi2023skincancer} and eMoglobin's machine learning pipeline for anemia~\cite{zhao2024emoglobin}. The 2022--2024 period saw CNN-ensemble and transfer learning approaches, such as STMap+CNN for SpO$_2$~\cite{cheng2024contactless}. From 2024 onward, Vision Transformers (ViT, T2T-ViT) have gained prominence, particularly for temporal (rPPG signal) and video (facial expression, gait) analysis.

However, lightweight CNN models (MobileNetV2, EfficientNet) retain advantages for on-device inference---MoodCapture~\cite{nepal2024moodcapture} uses EfficientNet specifically for privacy-preserving edge deployment. We anticipate hybrid architectures (CNN feature extraction + Transformer sequence modeling) becoming the dominant paradigm for smartphone-based health applications.

Three persistent technical barriers limit clinical translation:

\textbf{Illumination sensitivity}: Acharya et al.~\cite{acharya2025reliability} demonstrated significant rPPG reliability degradation under low illumination and elevated heart rate, conditions common in home monitoring scenarios. Adaptive illumination compensation and multi-color-space processing (YCrCb, CIELAB) are necessary but insufficient solutions.

\textbf{Skin tone bias}: FibriCheck~\cite{fibricheck2025} reported reduced AF detection sensitivity in darker skin tones, a bias rooted in training data predominantly representing lighter Fitzpatrick types. This systemic issue affects all camera-based biomarkers and demands Fitzpatrick-stratified evaluation as a minimum reporting standard.

\textbf{Motion artifacts}: Head movement and facial expressions introduce noise into rPPG signals. While benchmark datasets such as PURE include six motion conditions, real-world variability far exceeds laboratory protocols. Multi-ROI tracking and face stabilization algorithms partially address this challenge.

\subsection{Clinical Translation Gap}

A striking finding is the disconnect between reported accuracy and clinical readiness. Of 33 identified biomarkers, only two have achieved regulatory clearance: FibriCheck for AF~\cite{fibricheck2025} and select diabetic retinopathy screening systems. Several factors contribute to this gap:

\textbf{Sample size limitations}: The median validation cohort across Tier~2 and Tier~3 studies is approximately 200 participants, far below the thousands typically required for regulatory submissions. Only four studies exceed $n=500$.

\textbf{Laboratory vs.\ real-world performance}: Most high-accuracy results (e.g., stress detection at 95.83\%~\cite{fontes2024stress}, emotion classification across 16 categories~\cite{facialemotion2025}) were achieved on curated benchmark datasets. MoodCapture's~\cite{nepal2024moodcapture} in-the-wild accuracy of 60\% illustrates the performance degradation in uncontrolled conditions.

\textbf{Longitudinal evidence}: Nearly all studies employ cross-sectional designs. MoodCapture's 90-day longitudinal collection is the longest in our review. Chronic conditions such as Parkinson's disease, diabetes, and mental health disorders fundamentally require longitudinal monitoring data for clinical utility.

\textbf{Cost-effectiveness}: No study in our review conducted a formal health-economic analysis of camera-based screening versus standard care, a prerequisite for healthcare system adoption.

\subsection{Women's Health Opportunity: Endometriosis and PCOS}

Our systematic review revealed a complete absence of smartphone camera-based biomarker research targeting endometriosis or PCOS---conditions affecting approximately 10\% and 8--13\% of reproductive-age women, respectively. This gap is particularly notable given strong physiological rationale for camera-based monitoring:

\textbf{rPPG HRV for autonomic dysregulation}: PCOS patients exhibit significantly reduced SDNN, RMSSD, and HF power indicating sympathetic dominance, while endometriosis patients show vagal-mediated HRV reduction correlating with pelvic pain intensity. Given that rPPG-based HRV extraction achieves $r=0.85$--$0.95$ against ECG~\cite{rppgreview2024frontiers}, a daily 30-second rPPG selfie synchronized with menstrual cycle tracking could capture autonomic nervous system signatures of anovulation (PCOS) and inflammatory flares (endometriosis) without wearable devices.

\textbf{Facial skin analysis for hormonal fluctuation}: PCOS manifests with androgen-excess phenotypes including acne, hyperpigmentation, and hirsutism. Periodic standardized selfies analyzed by CNNs could track temporal skin changes as proxy indicators of hormonal fluctuation---leveraging established skin analysis technology~\cite{naqvi2023skincancer} for a novel clinical application.

\textbf{Nocturnal SpO$_2$ for PCOS-OSA comorbidity}: The elevated OSA prevalence in PCOS, combined with rPPG SpO$_2$ estimation capabilities~\cite{cheng2024contactless} and craniofacial OSA screening~\cite{osa2025jcsm}, suggests a camera-based nocturnal monitoring paradigm for metabolic risk stratification.

We propose a phased research agenda: (1) a 6--12 month pilot cohort integrating rPPG HRV with menstrual cycle app data (highest feasibility), (2) an exploratory facial skin analysis study for PCOS hormonal tracking (12--18 months), (3) nocturnal SpO$_2$ monitoring for PCOS-OSA screening (12--18 months), and (4) a multimodal integration model combining camera, app, and wearable data (18--24 months).

\subsection{Limitations}

This review has several limitations. First, the search was conducted by a single research team, and while we implemented a hallucination guard protocol, some partially verified references retain minor metric uncertainties. Second, our five-dimensional scoring framework, while systematic, involves subjective judgment, particularly for regulatory readiness assessment. Third, our search period (2015--2026) may miss foundational signal-processing work predating deep learning. Fourth, the women's health research gap analysis is based on physiological rationale rather than empirical evidence, and the proposed biomarkers require rigorous clinical validation. Finally, publication bias toward positive results likely inflates reported performance metrics across the reviewed literature.

%==============================================================================
\section{Conclusion}
%==============================================================================

This paper presents the first comprehensive systematic review and tiered evaluation of smartphone camera-based digital biomarkers for disease prediction, analyzing 40 verified studies across six technology domains. Our five-dimensional evaluation framework identified five Tier~1 biomarkers ready for clinical deployment, 13 Tier~2 biomarkers requiring multi-center validation, and 11 Tier~3 exploratory markers.

The field demonstrates remarkable technological progress---from early signal-processing rPPG to Transformer-based architectures achieving clinical-grade accuracy for heart rate, atrial fibrillation, blood oxygen saturation, neonatal jaundice, and diabetic retinopathy screening. However, a significant translation gap persists between laboratory performance and real-world clinical utility, with only two biomarkers achieving regulatory clearance.

We identify women's health---specifically endometriosis and PCOS---as the most critical unexplored frontier, where the convergence of rPPG-based autonomic monitoring, facial skin analysis, and menstrual cycle integration could enable non-invasive disease prediction using nothing more than a smartphone camera. Our proposed four-phase research roadmap offers a concrete pathway from this review to clinical impact.

As smartphone cameras continue to improve in resolution, frame rate, and computational capability, and as edge AI enables on-device privacy-preserving inference, we anticipate camera-based digital biomarkers becoming integral to preventive healthcare within the coming decade.

\bibliographystyle{IEEEtran}
\bibliography{references}

\end{document}
